% QMB_n05_bell_ffgft.tex — Kapitel 5: FFGFT-Antwort auf Bell
% Quellen: Dok. 230, 190, 353, 364
\chapter{Die FFGFT-Antwort auf Bell}
\label{ch:bell_ffgft}

\section{Topologische Verbindung statt Fernwirkung}

Die heutige FFGFT-Position (Dok.~230, Mai 2026) greift keine der drei
Bell-Annahmen an. Sie verändert stattdessen den ontologischen Rahmen:
Verschränkte Teilchen haben keinen getrennten Träger, sondern teilen
denselben Torus $T^4/\mathbb{Z}_3$. Die Korrelation ist eine
topologische Eigenschaft dieses gemeinsamen Trägers.

Das bedeutet konkret:
\begin{itemize}
  \item \textbf{Lokalität} gilt im Sinne des Trägers: Kein Signal
        überträgt Information schneller als Licht.
  \item \textbf{Realismus} gilt in dem Sinn, dass die Moden-
        konfiguration vor der Messung existiert.
  \item \textbf{Messfreiheit} wird nicht eingeschränkt.
\end{itemize}
Was aufgegeben wird, ist die Vorstellung, dass jedes Teilchen
einen \emph{unabhängigen} Träger hätte. Die topologische Verbindung
ist keine Fernwirkung — sie ist eine geometrische Tatsache.\status{B}

\section{Statistische Vorhersagen}

Die FFGFT reproduziert die Born-Regel und damit alle statistischen
Vorhersagen der QM vollständig:\status{B}
\begin{equation}
  P_{\mathrm{FFGFT}}(z=+1) = \tfrac{1+z}{2} = |\alpha|^2 = P_{\mathrm{QM}}(0).
\end{equation}
Insbesondere wird die CHSH-Schranke $2\sqrt2$ erreicht — nicht
unterschritten, wie es die frühen Dämpfungsansätze vorhergesagt
hatten. Der Bell-Test-Streit über lokalen Realismus löst sich
geometrisch auf: Es gibt keine Nichtlokalität, weil es keinen
getrennten Träger gibt.

\section{Die $\xi$-Abweichung}

Eine kleine Abweichung von der Tsirelson-Schranke ist aus der
$\theta_4$-Projektionsgeometrie des Messacts ableitbar. Der
Kandidat für die elementare Abweichung pro Messung ist:\status{S}
\begin{equation}
  \Delta\mathrm{CHSH}_{\mathrm{elem}} = \frac{\xi}{2\pi}
  \approx 2{,}1\times10^{-5}.
\end{equation}
Das ist die Phasenkosten eines $\theta_4$-Zyklus in geometrischen
Einheiten (Dok.~190, Präzisierung~12). Kumulativ über $N$ Messläufe
wächst diese Abweichung als:
\begin{equation}
  \Delta\mathrm{CHSH}(N) \approx \xi\,\frac{\ln N}{D_f}\cdot2\sqrt2.
\end{equation}
Bei $N=73$ ergibt das $\approx5\times10^{-4}$ — eine andere
Observable als die elementare, beide unter dem NISQ-Rauschen
($\sim10^{-2}$).

\textbf{Experimentelle Reichweite:} Hardware-Messungen (IBM Heron~r2,
Mai 2026, Dok.~147) zeigen Bell-Verletzung $S=2{,}74$ ($96{,}9\,\%$
der Tsirelson-Schranke). Der $\xi$-Effekt selbst erfordert
sub-ppm-stabile fehlerkorrigierte Qubits — nicht mit heutiger
NISQ-Hardware auflösbar.\status{K}

\begin{laierspur}
Die FFGFT sagt: Die starken Quantenkorrelationen sind echt und
korrekt — die QM hat recht. Was sie anders sieht: Die Korrelation
kommt nicht durch eine Verbindung \emph{zwischen} den Teilchen,
sondern weil die Teilchen nie wirklich getrennt waren. Sie sitzen
auf demselben geometrischen Träger wie zwei Punkte auf derselben
Trommel.
\end{laierspur}
\quelle{Dok.~190 (Präz.~11, 12), 230, 353, 364, 365}
